\subsection{全书符号与张量形状速览}
\phantomsection\label{app:symbols}
从本页开始，主线统一\textbf{每个 Token 占一行}。
单条输入写作 $\vect x^\top$，整段输入为 $X$；
投影写作 $\vect x^\top W$ 或 $XW$。
仅在点积、求梯度或显式说明转置时，用 $\vect x$ 表示同一向量的列形式。
\hyperref[intro:04-key-value-outer-product]{联想记忆专题}也使用行向量：写入 $\mathcal{M}=\sum_i k_i^\top v_i$，读出 $x\mathcal{M}$。

\begin{center}\small
\renewcommand{\arraystretch}{1.25}
\tagpdfsetup{table/header-rows={1}}
\begin{tabularx}{\linewidth}{@{}p{27mm}p{39mm}X@{}}
\toprule
符号 & 含义 & 典型形状\\\midrule
$T$ & 当前序列 Token 数 & $X:T\times d_{\rm model}$\\
$N$ & Query 位置数 & $Q:N\times d_k$\\
$M$ & Key/Value 位置数 & $K:M\times d_k$，$V:M\times d_v$\\
$d_{\rm model}$ & 每个 Token 的隐藏宽度 & $W_Q:d_{\rm model}\times d_k$\\
$d_k$ & Query 与 Key 的匹配宽度 & $QK^\top:N\times M$\\
$d_v$ & Value 的特征数 & $AV:N\times d_v$\\
$d_{\rm ff}$ & FFN 内部宽度 & $W_{\rm up}:d_{\rm model}\times d_{\rm ff}$\\
$|\mathcal V|$ & 词表大小 & $W_{\rm out}:d_{\rm model}\times|\mathcal V|$\\
\bottomrule
\end{tabularx}
\end{center}
Self-Attention 中 $N=M=T$；跨序列读取时 $N$ 与 $M$ 可以不同。
$\softmax$ 对矩阵默认逐行归一化。
温度 $\tau>0$ 与逆温度 $\beta=1/\tau$ 会在采样章继续出现。

\subsubsection*{同一个符号，只承担一个主角色}
\begin{center}\small
\renewcommand{\arraystretch}{1.2}
\tagpdfsetup{table/header-rows={1}}
\begin{tabularx}{\linewidth}{@{}p{32mm}X@{}}\toprule
符号 & 全书约定\\\midrule
$W_{\rm out},b_{\rm out}$ & 词表输出头；$W_O$ 专指多头 Attention 的输出投影\\
$Z,\vect z$ & logits；$L$ 专指损失，层数写 $N_{\rm layer}$\\
$P_{\rm pos}$ & 位置向量表；第 $i$ 行写 ${P_{\rm pos}}_{i,:}$\\
$\Pi,p_j$ & 整张词表概率矩阵、单行中的第 $j$ 类概率\\
$S,A$ & 注意力匹配分数、逐行 Softmax 后的读取权重\\
$\mathcal{M}$ & 外积记忆／线性注意力的状态，与分数 $S$ 区分\\
$F^{(\ell)},F_{\rm mid}$ & MLP 的第 $\ell$ 层激活、FFN 中间表示\\
$\mathcal A_i$ & 奖励归一化后的相对优势，与注意力权重 $A$ 区分\\
\bottomrule\end{tabularx}
\end{center}
脚本为保持接口简短，使用 \texttt{P} 对应 $P_{\rm pos}$、
\texttt{probs} 对应 $\Pi$、\texttt{Wout} 对应 $W_{\rm out}$；
大小写相近的 \texttt{Wo} 对应 $W_O$。代码名与数学名在此一一对照。

\subsubsection*{全书色彩表：用标签和线型共同识别}
\phantomsection\label{app:color-semantics}
\begin{center}\small
\tagpdfsetup{table/header-rows={1}}
\begin{tabularx}{\linewidth}{@{}p{26mm}X@{}}\toprule
类别 & 图色与冗余标记\\\midrule
角色 & \textcolor{blue}{Q 蓝}、\textcolor{teal}{K 青}、\textcolor{orange}{V 橙}，始终印角色字母\\
可见性 & \textcolor{visiblecolor}{可见绿／勾号}；\textcolor{maskedcolor}{遮蔽灰／叉号}\\
数据流 & \textcolor{flowcolor}{主流蓝灰／实箭头}；\textcolor{residualcolor}{残差紫／旁路}\\
数据与窗口 & \textcolor{datacolor}{数据玫红／点或柱}，目标柱另加粗边；噪声用黑色点；窗口用橄榄色实框／虚框\\
\bottomrule\end{tabularx}
\end{center}
背景的浅灰只分组，不承载角色语义。灰度阅读时按字母、叉／勾、边框与线型判断。

\begin{warningbox}
$QK^\top$ 是匹配分数，$\softmax(QK^\top)V$ 才是读取结果。
矩阵乘法顺序不能任意交换。
\end{warningbox}



